\documentclass[11pt]{article}
\usepackage[a4paper,margin=27mm]{geometry}
\usepackage{amsmath,amssymb,amsthm,hyperref}
\newtheorem{lemma}{Lemma}
\newtheorem{theorem}{Theorem}
\title{The equality case in the weak triple-minimum support bound}
\author{Research note: a general Gram-matrix obstruction}
\date{30 September 2026}
\begin{document}
\maketitle
This is the earlier uniform-weight argument, retained as an independent
arithmetic proof. The stronger theorem for arbitrary positive weights
is now proved in \path{weak_minimum_equality_obstruction.tex}.
Consider $q=3h$ equally weighted permutations on $n=q+1$ labels such
that each label is earliest in every triple containing it in exactly
$h$ rows. Write $r_{ij}$ for the number of rows with $i<j$.
Necessarily $h\le r_{ij}\le2h$.
Call a label bad if $r_{ij}=h$ for some opponent $j$, and good otherwise.
A bad label is globally first in exactly $h$ rows: the event $i<j$
already has the same size as every event $i<j,i<k$. Thus there are at
most three bad labels, and at least $q-2$ good labels.

\begin{lemma}\label{gram}
For a good pivot $i$, put $e_j=r_{ij}-h\in\{1,\ldots,h\}$ over its
$q$ opponents. Then
\begin{equation}\label{identity}
 q=\sum_j e_j+h-
       \frac{h(q-1)^2}{1+h\sum_j e_j^{-1}},
\end{equation}
and the positive integer
\begin{equation}\label{square}
 \left(\prod_j e_j\right)\left(1+h\sum_j e_j^{-1}\right)
\end{equation}
is a square.
\end{lemma}
\begin{proof}
Let $C$ be the $q\times q$ zero--one matrix whose row is a permutation
and whose opponent column $j$ indicates $i<j$. Triple-minimum
uniformity gives
\[
 C^{\mathsf T}C=D+h\mathbf1\mathbf1^{\mathsf T}=G,
 \qquad D=\operatorname{diag}(e_j),
 \qquad C^{\mathsf T}\mathbf1=e+h\mathbf1.
\]
Since $D$ is positive definite, $C$ is invertible. Consequently
\[
 q=(e+h\mathbf1)^{\mathsf T}G^{-1}(e+h\mathbf1).
\]
Sherman--Morrison gives \eqref{identity}; the matrix determinant lemma
gives \eqref{square} as $\det G=(\det C)^2$.
\end{proof}

\begin{lemma}\label{lasttwo}
Every good label avoids the last two positions in every row.
Consequently there are exactly three bad labels.
\end{lemma}
\begin{proof}
For the matrix $C$ of Lemma~\ref{gram}, solve $Cw=\mathbf1$.
Sherman--Morrison applied to
$Gw=C^{\mathsf T}\mathbf1=e+h\mathbf1$ gives
\[
 w_j=1-\frac{\alpha}{e_j},\qquad
 \alpha=\frac{h(q-1)}{1+h\sum_j e_j^{-1}}>0.
\]
Thus every $w_j<1$. A row of $C$ with zero or one entries equal to
one cannot have scalar product one with $w$. The pivot therefore
has at least two opponents after it in every row.

The last two positions must both be bad labels. If there were at
most two bad labels, they would occupy these positions in every
row and could never be first. This contradicts the fact that each
bad label is first in $h>0$ rows. There are at most three bad labels,
so there are exactly three.
\end{proof}

\begin{lemma}\label{threebad}
If there are exactly three bad labels, then every good pivot has
exactly three entries $e_j=h$.
\end{lemma}
\begin{proof}
The bad labels account for all first positions. For any good label $a$,
sum its minimum counts over the three pairs of bad labels. The total
is $3h=q$. In any row at least one bad label precedes $a$, so that row
contributes at most one to the sum. Every row must therefore
contribute one: $a$ follows the first bad label and precedes the other
two. Each bad label is first $h$ times, hence $r_{ab}=2h$ for every
good $a$ and bad $b$. Conversely, $e_j=h$ means $r_{ji}=h$, so $j$
must be bad.
\end{proof}

\begin{theorem}
For every integer $h\ge2$, no weak triple-minimum family of $3h$
equally weighted permutations exists on $3h+1$ labels.
\end{theorem}
\begin{proof}
By Lemmas~\ref{lasttwo} and~\ref{threebad}, there are exactly three
bad labels. Each row consists of its first bad label, all good
labels, and the remaining two bad labels. Partition the rows into
three groups of size $h$ according to their first bad label.

Fix two distinct good labels $a,c$, and write $r=r_{ac}$. For each
bad label $b$, let $x_b$ count rows with $a<c$ in the group beginning
with $b$. In all other rows both $a,c$ precede $b$. The minimum
condition on the triple $\{a,c,b\}$ gives $r-x_b=h$. Summing over
the three groups gives $r=\sum_bx_b=3(r-h)$, and hence
\[
 r=3h/2,\qquad e_{ac}=h/2.
\]
There are $3h-2\ge4$ good labels, so such a pair exists and $h$
must be even.

Fix a good pivot. Its three bad opponents have $e_j=h$, while its
other $3h-3$ opponents have $e_j=h/2$. Formula~\eqref{square} now
gives
\[
 (\det C)^2=\det G
 =h^3(h/2)^{3h-3}\bigl(1+3+2(3h-3)\bigr)
 =\frac{h^{3h}(3h-1)}{2^{3h-4}}.
\]
Because $h$ is even, both $h^{3h}$ and $2^{3h-4}$ are squares.
Thus $3h-1$ must be a rational square, and therefore an integer
square. This is impossible modulo three.
\end{proof}

\paragraph{Support consequence.}
Every equally weighted weak triple-minimum family on $n\ge5$ labels
has at least $n$ rows. Indeed, the general rank argument in
\path{../partial_fairness/weak_triple_minimum_support.tex} gives
$q\ge n-1$ even for unequal positive weights. In the equal-weight
case $3\mid q$, and equality $q=n-1$ would fall under the theorem.
The exception $h=1$ is real: on four labels the three permutations
$1203,3201,0213$ are weak triple-minimum uniform.

The new equality obstruction uses equal row weights, integrality of
the order counts, and the square determinant of an integer matrix.
It makes no stronger claim about unequal positive row weights.
The earlier note \path{six_weak_minimum_permutations.tex} gives an
independent position-based proof of the $h=2$ exclusion. The exact
small-profile enumeration in \path{check_weak_minimum_profiles.py}
and its JSON output is retained as a computational companion; the
general proof above does not depend on it.
\end{document}
